%--------------------------------------------------------------------------------
%
%
%--------------------------------------------------------------------------------
\pagebreak
\unnumberedSection{STEP - Step a T64 module}

%--------------------------------------------------------------------------------
\begin{iPageDescEntry}{Syntax}
\texttt{STEP [ <steps> [ , <modNum> ]} \\
\end{iPageDescEntry}

%--------------------------------------------------------------------------------
\begin{iPageDescEntry}{Description}

The \texttt{STEP} command advances a halted T64 module by executing one or more 
instructions on a processor module or step units on an I/O module. There can be 
one or more processors and I/O modules configured. When the step command is 
invoked without an argument, the simulator executes exactly one instruction on
all modules and return s control to the simulator. Likewise, a step count greater
than one will execute that amount of steps. 

If the module number is provided and the module is halted, only the module specified 
will execute the entered steps.

In either case, after the requested number of steps have completed, execution 
halts and control returns to the simulator command interface, allowing inspection 
of state, registers, memory, or windows. After inspection the step command or 
the run command resumes execution of the processor modules.

\end{iPageDescEntry}

%--------------------------------------------------------------------------------
\begin{iPageDescOpEntry}{Example}
\begin{lstlisting}[style=iPageOpStyle]

(5)->s        # execute a single instruction on all modules.
(6)->

(6)->s 10     # execute ten instructions in all modules.
(7)->

(8)->s 10, 5  # execute ten instructions on module 5
(9)->

\end{lstlisting}
\end{iPageDescOpEntry}

%--------------------------------------------------------------------------------
\begin{iPageDescEntry}{Notes}

The \texttt{STEP} command respects breakpoint and trace settings. A breakpoint 
or trace exception encountered during stepping may cause execution to stop before
the requested number of steps has completed.

\end{iPageDescEntry}
